\documentclass[11pt,a4paper]{article}

%================================================
% ENCODAGE ET LANGUE
%================================================

\usepackage[utf8]{inputenc}
\usepackage[T1]{fontenc}
\usepackage[french]{babel}

%================================================
% MISE EN PAGE
%================================================

\usepackage{geometry}
\geometry{margin=1.55cm}

%================================================
% PACKAGE GTOLI
%================================================

\usepackage{../styles/gtoli}

%================================================
% INFORMATIONS DU DOCUMENT
%================================================

\title{\textbf{Cours -- 5e secondaire}\\[0.2cm]
\large Titre du chapitre}

\author{GToli}
\date{\today}

%================================================
% DOCUMENT
%================================================

\begin{document}

\maketitle

%================================================
% OBJECTIFS
%================================================

\begin{cadreobjectif}
À la fin de ce chapitre, l’élève doit être capable de :
\begin{itemize}
    \item mobiliser les notions théoriques avec précision ;
    \item appliquer un protocole de résolution structuré ;
    \item justifier les choix méthodologiques ;
    \item interpréter les résultats obtenus ;
    \item résoudre des exercices de transfert ;
    \item rédiger une solution claire, rigoureuse et argumentée.
\end{itemize}
\end{cadreobjectif}

%================================================
\section{Notion concernée et enjeux du chapitre}
%================================================

\begin{cadrebleu}{Mise en contexte}
Cette section introduit la notion étudiée, ses enjeux et son rôle dans la progression du cours.

\medskip

Elle permet notamment de préciser :
\begin{itemize}
    \item la place du chapitre dans le programme ;
    \item les liens avec les notions déjà vues ;
    \item les compétences mobilisées ;
    \item les difficultés attendues.
\end{itemize}
\end{cadrebleu}

%================================================
\section{Prérequis}
%================================================

\begin{cadregris}{Notions à maîtriser avant de commencer}
Avant d’aborder ce chapitre, il est nécessaire de maîtriser :
\begin{itemize}
    \item les définitions fondamentales ;
    \item les propriétés déjà rencontrées ;
    \item les règles de calcul utiles ;
    \item les conventions de notation ;
    \item les méthodes élémentaires liées au chapitre.
\end{itemize}
\end{cadregris}

%================================================
\section{Rappel théorique structuré}
%================================================

\begin{cadretheorie}
Cette section contient les éléments théoriques indispensables :

\begin{itemize}
    \item définitions ;
    \item propriétés ;
    \item théorèmes ;
    \item formules ;
    \item conditions d’application ;
    \item conventions ;
    \item liens entre les notions.
\end{itemize}

\medskip

\textbf{Point important :} en 5e secondaire, la théorie doit être mobilisée non seulement pour calculer, mais aussi pour justifier les étapes du raisonnement.
\end{cadretheorie}

%================================================
\section{Méthode générale de résolution}
%================================================

\begin{cadremethode}
Pour résoudre un exercice de ce chapitre, il convient de suivre le protocole suivant :

\begin{enumerate}
    \item analyser précisément la consigne ;
    \item identifier les objets mathématiques ou scientifiques en jeu ;
    \item repérer les données explicites et implicites ;
    \item choisir la propriété ou la méthode adaptée ;
    \item vérifier les conditions d’application ;
    \item effectuer les calculs ou le raisonnement ;
    \item contrôler la cohérence du résultat ;
    \item rédiger une conclusion interprétée.
\end{enumerate}
\end{cadremethode}

%================================================
\section{Tableau de synthèse des notions}
%================================================

\begin{cadregris}{Tableau à compléter ou à adapter}
\renewcommand{\arraystretch}{1.35}
\begin{tabularx}{\textwidth}{>{\bfseries}p{3.5cm}X X}
\toprule
Notion & Signification & Méthode associée \\
\midrule
Notion 1 & Définition ou interprétation & Méthode à appliquer \\
Notion 2 & Définition ou interprétation & Méthode à appliquer \\
Notion 3 & Définition ou interprétation & Méthode à appliquer \\
\bottomrule
\end{tabularx}
\end{cadregris}

%================================================
\section{Exemple modèle}
%================================================

\begin{cadreexemple}
\textbf{Énoncé :}

Rédiger ici un exemple représentatif du chapitre.

\medskip

\textbf{Analyse de l’énoncé :}
\begin{itemize}
    \item Que demande la question ?
    \item Quelles sont les données disponibles ?
    \item Quelle notion du cours doit être mobilisée ?
    \item Quelle méthode est la plus pertinente ?
\end{itemize}

\medskip

\textbf{Résolution commentée :}

La résolution doit être présentée étape par étape en explicitant les choix effectués.
\end{cadreexemple}

%================================================
\section{Justification et rédaction}
%================================================

\begin{cadrevert}{Exigence de rigueur}
En 5e secondaire, une réponse correcte ne se limite pas au résultat final.

Une solution complète doit contenir :
\begin{itemize}
    \item une identification claire de la méthode ;
    \item des calculs organisés ;
    \item des justifications ;
    \item une interprétation éventuelle ;
    \item une conclusion rédigée.
\end{itemize}
\end{cadrevert}

%================================================
\section{Erreurs fréquentes}
%================================================

\begin{cadreattention}
Les erreurs fréquentes doivent être identifiées explicitement :

\begin{itemize}
    \item confusion entre deux notions proches ;
    \item oubli d’une condition d’existence ;
    \item application mécanique d’une formule ;
    \item absence de justification ;
    \item conclusion non interprétée ;
    \item erreur de signe, d’unité ou de notation.
\end{itemize}
\end{cadreattention}

%================================================
\section{Exercice d’application directe}
%================================================

\begin{cadreenonce}
\textbf{Exercice 1 :}

Rédiger ici un exercice d’application directe de la méthode.
\end{cadreenonce}

\begin{cadreaide}
Pour commencer :
\begin{itemize}
    \item repérer la notion principale ;
    \item identifier les données utiles ;
    \item choisir la formule ou la propriété ;
    \item vérifier les conditions d’application.
\end{itemize}
\end{cadreaide}

%================================================
\section{Solution détaillée}
%================================================

\begin{cadresolution}
La solution détaillée doit présenter :
\begin{enumerate}
    \item l’analyse de l’énoncé ;
    \item le choix de la méthode ;
    \item les calculs ;
    \item les justifications ;
    \item la conclusion.
\end{enumerate}
\end{cadresolution}

%================================================
\section{Exercices progressifs}
%================================================

\begin{cadreactivite}
\begin{enumerate}
    \item Exercice d’application directe.
    \item Exercice avec adaptation de la méthode.
    \item Exercice nécessitant plusieurs étapes.
    \item Exercice de transfert.
    \item Exercice de synthèse.
    \item Exercice de préparation à l’évaluation.
\end{enumerate}
\end{cadreactivite}

%================================================
\section{Exercice de transfert}
%================================================

\begin{cadreenonce}
\textbf{Exercice de transfert :}

Cet exercice doit obliger l’élève à adapter la méthode et non à simplement reproduire l’exemple.
\end{cadreenonce}

\begin{cadreaide}
Questions pour guider le raisonnement :
\begin{itemize}
    \item En quoi cet exercice diffère-t-il de l’exemple ?
    \item Quelle partie de la méthode reste valable ?
    \item Quelle adaptation faut-il effectuer ?
    \item Comment vérifier la cohérence de la réponse ?
\end{itemize}
\end{cadreaide}

%================================================
\section{Synthèse finale}
%================================================

\begin{cadresynthese}
À retenir pour ce chapitre :

\begin{itemize}
    \item la notion centrale est : \dotfill
    \item la méthode principale est : \dotfill
    \item les conditions d’application sont : \dotfill
    \item l’erreur la plus fréquente est : \dotfill
    \item la stratégie de vérification est : \dotfill
\end{itemize}
\end{cadresynthese}

%================================================
\section{Auto-évaluation}
%================================================

\begin{cadregris}{Je vérifie mon niveau de maîtrise}
\begin{itemize}
    \item[$\square$] Je connais les définitions importantes.
    \item[$\square$] Je sais choisir la méthode adaptée.
    \item[$\square$] Je sais justifier mes étapes.
    \item[$\square$] Je sais résoudre un exercice de base.
    \item[$\square$] Je sais résoudre une variante.
    \item[$\square$] Je sais interpréter mon résultat.
    \item[$\square$] Je sais rédiger une solution complète.
\end{itemize}
\end{cadregris}

%================================================
\section{Préparation à l’évaluation}
%================================================

\begin{cadreorange}{Conseils de révision}
Pour préparer l’évaluation :
\begin{enumerate}
    \item relire la théorie ;
    \item refaire l’exemple modèle sans regarder la solution ;
    \item refaire les exercices progressifs ;
    \item identifier les erreurs fréquentes ;
    \item s’entraîner sur un exercice de transfert ;
    \item rédiger les réponses avec justification.
\end{enumerate}
\end{cadreorange}

\end{document}